\documentclass[11pt,a4paper]{article}

\usepackage[T1]{fontenc}
\usepackage[utf8]{inputenc}
\usepackage[brazilian]{babel}
\usepackage{amsmath,amssymb,mathtools}
\usepackage{hyperref}
\usepackage[most]{tcolorbox}
\usepackage{geometry}
\usepackage{enumitem}
\usepackage{xcolor}
\geometry{margin=2.4cm}
\hypersetup{colorlinks=true,linkcolor=blue!60!black,urlcolor=blue!60!black}
\setlist[itemize]{topsep=4pt,itemsep=2pt}
\setlength{\parskip}{0.5em}
\setlength{\parindent}{0pt}
\newcommand{\dd}{\mathrm{d}}

\newtcolorbox{resultado}{
  colback=blue!4!white,
  colframe=blue!55!black,
  title=Resultado final,
  breakable
}

\newtcolorbox{intuicao}{
  colback=green!4!white,
  colframe=green!45!black,
  title=Intui\c{c}\~ao econ\^omica,
  breakable
}

\newtcolorbox{convencao}{
  colback=gray!8!white,
  colframe=gray!55!black,
  title=Conven\c{c}\~oes desta nota,
  breakable
}

\title{Cap\'itulo 1 --- Deriva\c{c}\~oes do Modelo de Solow}
\author{romer-study}
\date{\today}

\begin{document}

\maketitle
\tableofcontents
\bigskip

\section{Objetivo da nota}

Esta nota deriva, passo a passo, as equa\c{c}\~oes centrais do Cap\'itulo 1 do
Romer. O objetivo \'e sair da tecnologia agregada
\[
Y(t)=F\bigl(K(t),A(t)L(t)\bigr)
\]
e chegar \`a equa\c{c}\~ao em capital por trabalhador efetivo,
\[
\dot{k}(t)=s f\bigl(k(t)\bigr)-\bigl(n+g+\delta\bigr)k(t).
\]

Tamb\'em derivamos:
\begin{itemize}
  \item a forma intensiva \(y=f(k)\);
  \item o estado estacion\'ario;
  \item a solu\c{c}\~ao Cobb--Douglas;
  \item a Regra de Ouro;
  \item a decomposi\c{c}\~ao de growth accounting.
\end{itemize}

\begin{convencao}
A nota\c{c}\~ao coincide com a do reposit\'orio:
\[
\alpha,\quad s,\quad n,\quad g,\quad \delta.
\]
No c\'odigo, esses objetos aparecem na classe
\texttt{SolowModel} em \texttt{ch01\_solow/ch01\_solow.py}.
\end{convencao}

\section{Hip\'oteses b\'asicas}

Assumimos:
\begin{itemize}
  \item tecnologia com retornos constantes de escala;
  \item taxa de poupan\c{c}a ex\'ogena \(s\);
  \item crescimento populacional ex\'ogeno \(\dot L/L=n\);
  \item progresso t\'ecnico ex\'ogeno \(\dot A/A=g\);
  \item deprecia\c{c}\~ao \(\delta>0\).
\end{itemize}

O produto agregado \'e
\[
Y=F(K,AL),
\]
e a acumula\c{c}\~ao do capital agregado \'e
\[
\dot K=sY-\delta K.
\]

\section{Da fun\c{c}\~ao agregada \`a forma intensiva}

\subsection{Retornos constantes de escala}

Por hip\'otese, para todo \(\lambda>0\),
\[
F(\lambda K,\lambda AL)=\lambda F(K,AL).
\]
Escolhendo
\[
\lambda=\frac{1}{AL},
\]
obtemos
\[
F\!\left(\frac{K}{AL},1\right)=\frac{1}{AL}F(K,AL).
\]
Multiplicando ambos os lados por \(AL\),
\[
F(K,AL)=AL\cdot F\!\left(\frac{K}{AL},1\right).
\]

Agora definimos
\[
k\equiv \frac{K}{AL},
\qquad
y\equiv \frac{Y}{AL},
\qquad
f(k)\equiv F(k,1).
\]
Logo,
\[
Y=AL\cdot f(k)
\qquad \Longrightarrow \qquad
y=f(k).
\]

\begin{resultado}
Sob retornos constantes de escala,
\[
\frac{Y}{AL}=f\!\left(\frac{K}{AL}\right).
\]
Essa \'e a forma intensiva do modelo.
\end{resultado}

\begin{intuicao}
O modelo reduz a economia a uma rela\c{c}\~ao entre produto e capital por
trabalhador efetivo. O tamanho absoluto da economia deixa de importar.
\end{intuicao}

\section{Deriva\c{c}\~ao da din\^amica de \(k\)}

Come\c{c}amos com
\[
k=\frac{K}{AL}.
\]
Derivando no tempo,
\[
\dot{k}
=
\frac{\dot K}{AL}
-
\frac{K}{(AL)^2}\frac{\dd(AL)}{\dd t}.
\]

Como
\[
\frac{\dot A}{A}=g
\qquad \text{e} \qquad
\frac{\dot L}{L}=n,
\]
segue que
\[
\frac{\dd(AL)}{\dd t}=A\dot L+L\dot A=AL(n+g).
\]
Substituindo,
\[
\dot{k}
=
\frac{\dot K}{AL}
-
\frac{K}{AL}(n+g)
=
\frac{\dot K}{AL}-(n+g)k.
\]

Da lei de movimento de \(K\),
\[
\dot K=sY-\delta K.
\]
Dividindo por \(AL\),
\[
\frac{\dot K}{AL}
=
s\frac{Y}{AL}-\delta \frac{K}{AL}
=
s f(k)-\delta k.
\]
Portanto,
\[
\dot{k}
=
s f(k)-\delta k-(n+g)k
=
s f(k)-\bigl(n+g+\delta\bigr)k.
\]

\begin{resultado}
A equa\c{c}\~ao fundamental do modelo de Solow em unidades por trabalhador
efetivo \'e
\[
\dot{k}=s f(k)-\bigl(n+g+\delta\bigr)k.
\]
\end{resultado}

\section{Estado estacion\'ario}

Um steady state em termos de \(k\) satisfaz
\[
\dot{k}=0.
\]
Logo,
\[
0=s f(k^\ast)-\bigl(n+g+\delta\bigr)k^\ast,
\]
ou
\[
s f(k^\ast)=\bigl(n+g+\delta\bigr)k^\ast.
\]

\begin{resultado}
A condi\c{c}\~ao de estado estacion\'ario do Solow \'e
\[
s f(k^\ast)=\bigl(n+g+\delta\bigr)k^\ast.
\]
\end{resultado}

Uma vez obtido \(k^\ast\), temos
\[
y^\ast=f(k^\ast),
\qquad
i^\ast=s f(k^\ast),
\qquad
c^\ast=(1-s)f(k^\ast).
\]

\subsection{Crescimento balanceado}

Se \(k=K/(AL)\) \'e constante, ent\~ao \(K\) e \(AL\) crescem \`a mesma taxa:
\[
\frac{\dot K}{K}=n+g.
\]
Como \(Y=ALf(k)\) com \(k\) constante,
\[
\frac{\dot Y}{Y}=n+g.
\]
E o produto por trabalhador cresce \`a taxa do progresso t\'ecnico:
\[
\frac{\dot{(Y/L)}}{Y/L}=g.
\]

\section{Caso Cobb--Douglas}

Suponha
\[
Y=K^\alpha (AL)^{1-\alpha},
\qquad 0<\alpha<1.
\]
Dividindo por \(AL\),
\[
y
=
\frac{K^\alpha (AL)^{1-\alpha}}{AL}
=
\left(\frac{K}{AL}\right)^\alpha
=
k^\alpha.
\]
Assim,
\[
f(k)=k^\alpha.
\]

\subsection{Deriva\c{c}\~ao de \(k^\ast\)}

Partimos da condi\c{c}\~ao de steady state:
\[
s k^{\ast\alpha}=\bigl(n+g+\delta\bigr)k^\ast.
\]
Como \(k^\ast>0\), dividimos por \(k^\ast\):
\[
s k^{\ast(\alpha-1)}=n+g+\delta.
\]
Logo,
\[
k^\ast
=
\left(\frac{s}{n+g+\delta}\right)^{\frac{1}{1-\alpha}}.
\]

\subsection{Deriva\c{c}\~ao de \(y^\ast\) e \(c^\ast\)}

Como \(y^\ast=k^{\ast\alpha}\),
\[
y^\ast
=
\left(\frac{s}{n+g+\delta}\right)^{\frac{\alpha}{1-\alpha}}.
\]
E como \(c^\ast=(1-s)y^\ast\),
\[
c^\ast
=
(1-s)
\left(\frac{s}{n+g+\delta}\right)^{\frac{\alpha}{1-\alpha}}.
\]

\begin{resultado}
No caso Cobb--Douglas,
\[
k^\ast=\left(\frac{s}{n+g+\delta}\right)^{\frac{1}{1-\alpha}},
\qquad
y^\ast=\left(\frac{s}{n+g+\delta}\right)^{\frac{\alpha}{1-\alpha}},
\qquad
c^\ast=(1-s)y^\ast.
\]
\end{resultado}

\section{Regra de Ouro}

A Regra de Ouro escolhe o capital que maximiza o consumo de steady state:
\[
c(k)=f(k)-\bigl(n+g+\delta\bigr)k.
\]
A condi\c{c}\~ao de primeira ordem \'e
\[
\frac{\dd c(k)}{\dd k}
=
f'(k)-\bigl(n+g+\delta\bigr)=0.
\]
Portanto,
\[
f'(k_{\text{gold}})=n+g+\delta.
\]

\begin{resultado}
A condi\c{c}\~ao de Regra de Ouro \'e
\[
f'(k_{\text{gold}})=n+g+\delta.
\]
\end{resultado}

Se \(f(k)=k^\alpha\), ent\~ao
\[
f'(k)=\alpha k^{\alpha-1}.
\]
Logo,
\[
\alpha k_{\text{gold}}^{\alpha-1}=n+g+\delta
\qquad \Longrightarrow \qquad
k_{\text{gold}}=
\left(\frac{\alpha}{n+g+\delta}\right)^{\frac{1}{1-\alpha}}.
\]

Para obter a taxa de poupan\c{c}a correspondente, usamos
\[
s_{\text{gold}}f(k_{\text{gold}})
=
\bigl(n+g+\delta\bigr)k_{\text{gold}}.
\]
Substituindo a condi\c{c}\~ao marginal, segue que
\[
s_{\text{gold}}=\alpha.
\]

\begin{resultado}
No caso Cobb--Douglas,
\[
k_{\text{gold}}=
\left(\frac{\alpha}{n+g+\delta}\right)^{\frac{1}{1-\alpha}},
\qquad
s_{\text{gold}}=\alpha.
\]
\end{resultado}

\section{Contabilidade do crescimento}

Com Cobb--Douglas,
\[
Y=K^\alpha (AL)^{1-\alpha}.
\]
Tomando logaritmos,
\[
\ln Y = \alpha \ln K + (1-\alpha)\ln L + (1-\alpha)\ln A.
\]
Derivando no tempo,
\[
\frac{\dot Y}{Y}
=
\alpha \frac{\dot K}{K}
+
(1-\alpha)\frac{\dot L}{L}
+
(1-\alpha)\frac{\dot A}{A}.
\]

Na pr\'atica emp\'irica, o termo tecnol\'ogico costuma ser escrito como res\'iduo
de PTF:
\[
g_{\text{TFP}}
\equiv
g_Y-\alpha g_K-(1-\alpha)g_L.
\]
Logo,
\[
g_Y=\alpha g_K+(1-\alpha)g_L+g_{\text{TFP}}.
\]

\begin{resultado}
A decomposi\c{c}\~ao usada no projeto \'e
\[
g_{\text{TFP}}=g_Y-\alpha g_K-(1-\alpha)g_L.
\]
\end{resultado}

No reposit\'orio, isso aparece em
\texttt{SolowModel.growth\_accounting(...)}.

\section{Do livro para o c\'odigo}

As correspond\^encias principais s\~ao:
\begin{itemize}
  \item \(f(k)=k^\alpha\) \(\leftrightarrow\) \texttt{f(...)};
  \item \(f'(k)=\alpha k^{\alpha-1}\) \(\leftrightarrow\) \texttt{f\_prime(...)};
  \item \(\dot{k}=s f(k)-(n+g+\delta)k\) \(\leftrightarrow\) \texttt{k\_dot(...)};
  \item \(k^\ast,y^\ast,c^\ast\) \(\leftrightarrow\) \texttt{steady\_state(...)};
  \item \(k_{\text{gold}},s_{\text{gold}}\) \(\leftrightarrow\) \texttt{golden\_rule(...)}.
\end{itemize}

Em particular,
\[
\texttt{self.dilution}=n+g+\delta
\]
\'e a tradu\c{c}\~ao direta do termo de reposi\c{c}\~ao e dilui\c{c}\~ao.

\section{Fechamento}

O Cap\'itulo 1 do Romer concentra o crescimento de longo prazo numa \'unica
equa\c{c}\~ao para \(k\):
\[
\dot{k}=s f(k)-\bigl(n+g+\delta\bigr)k.
\]
Uma vez derivada essa express\~ao, seguem de forma organizada o steady state,
a Regra de Ouro, a converg\^encia e a contabilidade do crescimento.

\end{document}
